
We develop posterior sampling at the level of the general interpolant path of Section~\ref{subsec:flow_models}, so one construction covers stochastic interpolants, flow matching, and diffusion models, and yields both SDE and ODE samplers. The prior model is trained once and never retrained; it supplies a velocity $\fmvel_\tau$, score $\genscore_\tau$, or SI drift for the path onto $\conddist{\bx_1}{\bx_0}$, and we want samples from the one-step posterior $\conddist{\bx_1}{\obs,\bx_0}$ of Eq.~\eqref{eq:posterior_distribution}. The section mirrors the method: conditioning the path yields an exact family of posterior samplers driven by one likelihood-score term (\ref{subsec:posterior_dynamics}); that term is replaced by a closed-form observation-interpolant surrogate (\ref{subsec:obs_interpolation}); two covariance approximations make it cheap at field scale (\ref{subsec:approximations}); and \ref{subsec:summary_table} compares the resulting samplers.

\subsection{Posterior sampling: one guidance term, a family of samplers}
\label{subsec:posterior_dynamics}

Conditioning on $\obs$ preserves the interpolant path and only reweights it by the \emph{likelihood tilt}
\begin{align}\label{eq:likelihood_tilt}
    \Phi_\tau^{\obs}(\bx,\bx_0) := \conddist{\obs}{\bX_\tau=\bx,\bx_0}
    = \E\!\left[\conddist{\obs}{\bx_1}\,\middle|\,\bX_\tau=\bx,\bX_0=\bx_0\right],
\end{align}
the observation probability given the current state. The conditioned path is again an interpolant path with the same schedules and anchor, its target retargeted to the posterior, so its score and velocity are the prior ones shifted by the \emph{same} likelihood score $\nabla_{\bx}\log\Phi_\tau^{\obs}$ -- scaled by one for the score and by $\vscoef_\tau$ for the velocity (Proposition~\ref{prop:conditional_path}, Appendix~\ref{appendix:proof_conditioning}). This single term is all that conditioning changes, and the equal-marginal family \eqref{eq:general_sde_family} on the conditioned path becomes a family of posterior samplers:

\begin{theorem}[Unified posterior sampler]\label{theorem:unified_posterior}
Let the prior model define an interpolant path with velocity $\fmvel_\tau$ and score $\genscore_\tau$ whose terminal law is $\conddist{\bx_1}{\bx_0}$, and let $\Phi_\tau^{\obs}$ be the likelihood tilt \eqref{eq:likelihood_tilt}. For any diffusion schedule $\gdiff_\tau\ge0$, the SDE
\begin{align}\label{eq:unified_posterior_drift}
    \rd\bXstar_\tau
    = \underbrace{\left[\fmvel_\tau(\bXstar_\tau,\bx_0) + \tfrac12\gdiff_\tau^2\,\genscore_\tau(\bXstar_\tau,\bx_0)\right]}_{\text{prior drift }\drift^{\gdiff}_\tau}\rd\tau
    + \underbrace{\big(\vscoef_\tau + \tfrac12\gdiff_\tau^2\big)}_{=:\,\gweight_\tau}\nabla_{\bXstar_\tau}\log\Phi_\tau^{\obs}(\bXstar_\tau,\bx_0)\,\rd\tau
    + \gdiff_\tau\,\rd\bW_\tau ,
\end{align}
initialised from $\bXstar_0\sim p_0^{\obs}$ with $p_0^{\obs}(\rd\bx)\propto\Phi_0^{\obs}(\bx,\bx_0)\,p_0(\rd\bx)$, has marginals $p_\tau^{\obs}(\bx\mid\bx_0)\propto\Phi_\tau^{\obs}\,p_\tau$ for all $\tau\in[0,1]$, and in particular generates $\bXstar_1\sim\conddist{\bx_1}{\obs,\bx_0}$.
\end{theorem}

The dynamics split into the prior drift and one likelihood-score correction with \emph{guidance weight} $\gweight_\tau=\vscoef_\tau+\tfrac12\gdiff_\tau^2$: the velocity--score term plus the diffusion term. In the ideal setting (exact prior and tilt) \emph{every} member of the family yields exact posterior samples; the proof is in Appendix~\ref{appendix:proof_conditioning}, where the initial reweighting is also shown to be vacuous for both model classes. The prior drift, likelihood score, and observation handling are shared across the family; members differ only through $\gdiff_\tau$, and three are of practical interest:
\begin{align}\label{eq:three_drifts}
\begin{array}{l@{\;}l@{\quad}l}
\text{SI-SDE} &(\gdiff_\tau=\gamma_\tau): &
  \drift^{\gamma}_\tau\,\rd\tau
  +\ \big(\vscoef_\tau+\tfrac12\gamma_\tau^2\big)\,\nabla\log\Phi_\tau^{\obs}\,\rd\tau\ +\ \gamma_\tau\,\rd\bW_\tau, \\[2pt]
\text{DM-SDE} &(\gdiff_\tau>0): &
  \drift^{\gdiff}_\tau\,\rd\tau
  +\ \big(\vscoef_\tau+\tfrac12\gdiff_\tau^2\big)\,\nabla\log\Phi_\tau^{\obs}\,\rd\tau\ +\ \gdiff_\tau\,\rd\bW_\tau, \\[2pt]
\text{FM-ODE} &(\gdiff_\tau=0): &
  \fmvel_\tau\,\rd\tau
  +\ \vscoef_\tau\,\nabla\log\Phi_\tau^{\obs}\,\rd\tau .
\end{array}
\end{align}
At $\gdiff_\tau=0$ the noise and the prior-drift score term vanish, leaving the deterministic guided ODE; $\gdiff_\tau=\gamma_\tau$ recovers the native SI-SDE; any other $\gdiff_\tau>0$ on an FM prior is a guided denoising-diffusion sampler. Stochastic interpolants, denoising diffusion, and guided flows are one family, differing only in a scalar schedule.

\begin{remark}[Relation to Doob's $h$-transform]\label{rem:doob}
An alternative posterior SDE applies Doob's $h$-transform to a fixed prior diffusion, giving the correction weight $\gdiff_\tau^2$ instead of $\gweight_\tau$. It conditions trajectories rather than marginals -- its intermediate marginals are not $p_\tau^{\obs}$ -- and has no deterministic member; the marginal construction \eqref{eq:unified_posterior_drift} is also the one consistent with the likelihood surrogate below (Appendix~\ref{appendix:proof_posterior}).
\end{remark}

\subsection{A closed-form likelihood score from observation interpolation}
\label{subsec:obs_interpolation}

The correction in \eqref{eq:unified_posterior_drift} needs the likelihood score $\nabla_{\bx}\log\Phi_\tau^{\obs}$ at intermediate states, but the exact tilt
\begin{align}\label{eq:phi_integral}
    \Phi_{\tau}^{\obs}(\bx_\tau, \bx_0)
    = \int \conddist{\obs}{\bx_1}\, \conddist{\bx_1}{\bx_\tau, \bx_0}\, \rd \bx_1
\end{align}
is intractable: evaluating it requires integrating the dynamics forward and backpropagating through the trajectory. Our key construction is to interpolate the \emph{observation} exactly as the model interpolates the state, mirroring Eq.~\eqref{eq:general_interpolant} in observation space:
\begin{align} \label{eq:interpolated_observation}
    \interpolantobs_\tau(\bx_0, \obs) = \alpha_\tau \obsmatrix\bx_0 + \beta_\tau \obs,
    \qquad \interpolantobs_0 = \alpha_0\obsmatrix\bx_0,\quad \interpolantobs_1 = \obs,
\end{align}
which reduces to $\beta_\tau\obs$ for FM ($\alpha_\tau\equiv0$): at pseudo-time $\tau$, $\interpolantobs_\tau$ is the part of the observation the state should already explain. Replacing the tilt by the surrogate $\bar{\Phi}_{\tau}^{\interpolantobs_\tau}(\bx_\tau, \bx_0) = \conddist[\tau]{\interpolantobs_\tau}{\bx_\tau, \bx_0}$ gives the working form of \eqref{eq:unified_posterior_drift} for all three samplers,
\begin{align} \label{eq:interpolant_posterior_drift}
    \postdrift_\theta(\bx_{\tau}, \bx_0, \interpolantobs_\tau, \tau)
    =
    \drift^{\gdiff}_\tau(\bx_{\tau}, \bx_0)
    + \gweight_\tau\,
    \nabla_{\bx_\tau} \log\bar{\Phi}_{\tau}^{\interpolantobs_\tau}(\bx_\tau, \bx_0).
\end{align}
The surrogate is tractable because $\interpolantobs_\tau$ is a \emph{linear} function of the current state plus noise whose moments the trained model already knows. Writing the zero-mean source process as $\src_\tau := \sigpath_\tau\latentnoise$, so $\bx_\tau = \alpha_\tau\bx_0 + \beta_\tau\bx_1 + \src_\tau$:

\begin{lemma}[Interpolated-observation likelihood]\label{lem:interpolated_likelihood}
Let $\obsoperator$ be linear with matrix $\obsmatrix$ and $\obsnoise \sim \mathcal{N}(0,\obscov)$. Then
\begin{equation}\label{eq:ybar_decomp}
  \interpolantobs_\tau = \obsmatrix\bx_\tau + \interpolantobsnoise_\tau,
  \qquad
  \interpolantobsnoise_\tau := \beta_\tau\obsnoise - \obsmatrix\src_\tau,
\end{equation}
and the conditional mean and covariance of $\conddist{\interpolantobs_\tau}{\bx_0, \bx_\tau}$ are
\begin{align}
  \bar{\mu}_\tau(\bx_\tau, \bx_0)
  &= \obsmatrix\bx_\tau - \obsmatrix\,\srcmean_\tau(\bx_\tau,\bx_0),
     \label{eq:bias} \\[2pt]
  \bar{\Sigma}_{\tau}(\bx_\tau, \bx_0)
  &= \beta_\tau^2\obscov + \obsmatrix\,\srccov_\tau(\bx_\tau,\bx_0)\,\obsmatrix^\top,
     \label{eq:cov_correction}
\end{align}
with source moments $\srcmean_\tau := \E[\src_\tau\mid\bx_\tau,\bx_0]$ and $\srccov_\tau := \cov(\src_\tau\mid\bx_\tau,\bx_0)$.
\end{lemma}

\begin{lemma}[Source conditional moments]\label{lem:source_moments}
For the path \eqref{eq:general_interpolant}, the conditional moments of the source process are
\begin{align}\label{eq:source_moments}
    \srcmean_\tau(\bx,\bx_0) = -\sigpath_\tau^2\,\genscore_\tau(\bx,\bx_0),
    \qquad
    \srccov_\tau(\bx,\bx_0) = \sigpath_\tau^2\,\bI + \sigpath_\tau^4\,\nabla_{\bx}\genscore_\tau(\bx,\bx_0).
\end{align}
\end{lemma}

Both proofs are in Appendix~\ref{appendix:proof_interpolant_likelihood}; Lemma~\ref{lem:source_moments} is the first- and second-order Tweedie identity for the path, and holds \emph{identically} for every model class -- only the noise scale $\sigpath_\tau$ and how the score is read from the network (Eqs.~\eqref{eq:SI_score}, \eqref{eq:fm_score}) differ. Moreover $\bar\mu_\tau = \alpha_\tau\obsmatrix\bx_0 + \beta_\tau\obsmatrix\,\E[\bx_1\mid\bx_\tau,\bx_0]$: the likelihood mean is the observation interpolant at the model's denoiser, so \eqref{eq:interpolant_posterior_drift} is a Tweedie-style guidance term shared by all three samplers.

Taking the interpolant likelihood Gaussian with the moments of Lemma~\ref{lem:interpolated_likelihood}, the likelihood score is available in closed form,
\begin{align}\label{eq:interpolant_score_closed}
    \bar{S}_\tau
    = (\nabla_{\bx_\tau}\bar{\mu}_\tau)^\top\,\bar{\Sigma}_\tau^{-1}\,(\interpolantobs_\tau - \bar{\mu}_\tau),
    \qquad
    \bar{\Sigma}_\tau = \beta_\tau^2\,\obscov + \obsmatrix\,\srccov_\tau\,\obsmatrix^\top,
\end{align}
with the covariance held fixed in $\bx_\tau$ when differentiating, as is common in guidance methods. No inner integral, no trajectory backpropagation, and no tuned guidance strength: the pull towards the observation is set by the \emph{inflated} covariance $\bar\Sigma_\tau$, which augments the raw observation covariance $\beta_\tau^2\obscov$ by the model's source covariance. This is exactly the $\Pi$GDM-style inflation \cite{song_pseudoinverse_2023}, here free since the trained model already supplies $\srccov_\tau$; it carries the prior-covariance information that lets the guidance reach unobserved degrees of freedom. Diffusion posterior sampling (DPS) \cite{chung_diffusion_2023, chung_improving_2024} drops it, which costs little when observations are dense but loses substantial accuracy when they are sparse (Section~\ref{sec:results}).

\subsection{Making the guidance cheap: two covariance approximations}
\label{subsec:approximations}\label{subsec:covariance_approximation}

The only expensive object in \eqref{eq:interpolant_score_closed} is $\srccov_\tau$, whose second term is a network Jacobian. Two approximations, at opposite ends of the cost/accuracy axis, make the guidance practical; which to use is dictated by the observation regime.

\begin{corollary}[Jacobian-free covariance]\label{cor:cheap_drift}
Dropping the Jacobian term in \eqref{eq:source_moments}, $\srccov_\tau\approx\sigpath_\tau^2\bI$, replaces the inflated covariance by the state-independent, precomputable
\begin{align}\label{eq:cov_cheap}
    \bar{\Sigma}_\tau \approx \beta_\tau^2\,\obscov + \sigpath_\tau^2\,\obsmatrix\obsmatrix^\top,
\end{align}
which depends on $\tau$ only through the scalars $\beta_\tau^2,\sigpath_\tau^2$ and requires no network Jacobian.
\end{corollary}

The truncation is accurate when $\sigpath_\tau^2\|\nabla_{\bx}^2\log p_\tau\|\ll1$ (near $\tau\to1$, or near-Gaussian priors); with time-independent $\obsmatrix,\obscov$ the matrix $\obsmatrix\obsmatrix^\top$ is precomputed once and each step costs one cheap $\obsdim\times \obsdim$ solve. This is the method of choice for dense, informative observations, where the likelihood dominates and the samplers nearly coincide.

\paragraph{Ensemble-shared Jacobian.}
For sparse observations, the discarded Jacobian carries exactly the prior correlations the guidance needs, and Corollary~\ref{cor:cheap_drift} degrades. Retaining the full $\srccov_\tau$, however, costs $\obsdim$ Jacobian--vector products of the network \emph{per ensemble member} per step -- $\mathcal{O}(E\,\obsdim)$ network passes, prohibitive at field scale. We instead evaluate the source-covariance Jacobian \emph{once per step, at the ensemble mean}, and share the resulting $\obsdim\times\obsdim$ factorisation across all $E$ members: the cost drops to $\mathcal{O}(\obsdim)$ network passes and a single factorisation, and the approximation -- the analogue of the shared-gain analysis step of the EnKF -- is exact in the limit of a tight ensemble. It is what makes the accurate inflated covariance feasible in the sparse regime where it matters most (Section~\ref{sec:results}). Throughout, the guidance solves only $\obsdim\times\obsdim$ systems, never a state-space $\statedim\times\statedim$ one, so high-dimensional states pose no obstacle.

\subsection{Comparison of posterior samplers}
\label{subsec:summary_table}\label{subsec:supplying_diffusion}

The family \eqref{eq:interpolant_posterior_drift} contains three samplers of practical interest, sharing the observation-interpolant likelihood with inflated covariance and differing only in $\gdiff_\tau$, the weight $\gweight_\tau=\vscoef_\tau+\tfrac12\gdiff_\tau^2$, and the source. The SI-SDE ($\gdiff_\tau=\gamma_\tau$) reuses the native SDE \eqref{eq:SI_SDE}, whose drift already contains the score, from the point-mass source $\delta_{\bx_0}$. The DM-SDE ($\gdiff_\tau>0$) lifts the FM ODE by recovering the score via \eqref{eq:fm_score}, initialised exactly from $\mathcal{N}(0,\bI)$ since $\Phi_0^{\obs}$ is constant (Appendix~\ref{appendix:proof_conditioning}); an endpoint-vanishing $\gdiff_\tau\propto\sqrt{\sigpath_\tau\beta_\tau}$ keeps the recovered score finite as $\tau\to1$. The guided ODE ($\gdiff_\tau=0$) is the probability-flow ODE of the posterior path with weight $\vscoef_\tau$, coinciding with interpolant-guided flows for linear inverse problems \cite{yan_fig_2024, pokle_training-free_2024}; being deterministic it is well posed for FM but not for SI, whose point-mass source would collapse to a single point. Tables~\ref{tab:samplers}--\ref{tab:approximations} (Appendix~\ref{appendix:bookkeeping}) contrast the samplers, the model-specific bookkeeping, and the exactness regime of every approximation. In brief, the configuration is set by two orthogonal choices -- the sampler ($\gdiff_\tau$) and the covariance -- selected by problem geometry: dense observations favour the Jacobian-free covariance, sparse observations the ensemble-shared inflated one, the ODE is cheapest per step, and the SDE samplers give better-calibrated spread (decision guide in Figure~\ref{fig:decision_guide}, Appendix~\ref{appendix:bookkeeping}).
